\SourceLeaf{633}
\begin{center}[20]\par\medskip[\textit{C. F. Gauss an W. Weber.}]\end{center}
Hochgeschätzter Freund,

Seit Anfang dieses Jahres wird meine Zeit unablässig auf so vielfältige Weise beansprucht und zersplittert. Zugleich ist mein Gesundheitszustand für längeres Arbeiten so wenig günstig gewesen, dass ich bisher gar nicht dazu gekommen bin, den kleinen Aufsatz durchzugehen, den Sie mir vor zwei Monaten freundlicherweise geschickt haben. Erst jetzt konnte ich ihn flüchtig durchsehen. Dabei habe ich jedoch erkannt, dass das Thema zu denselben Untersuchungen gehört, mit denen ich mich vor etwa 10 Jahren (ich meine besonders 1834–1836) sehr eingehend beschäftigt habe. Um ein gründliches und erschöpfendes Urteil über Ihren Aufsatz abgeben zu können, genügt es nicht, \textit{diesen} durchzulesen. Ich müsste mich vielmehr erst wieder ganz in meine eigenen Arbeiten aus jener Zeit einarbeiten. Das würde umso länger dauern, als ich bei einer versuchsweisen Durchsicht meiner Papiere bisher nur einige Bruchstücke gefunden habe, obwohl wahrscheinlich noch viel mehr vorhanden ist, wenn auch nicht in vollständig geordneter Form.

Wenn ich mir aber, nachdem mir diese Gegenstände seit mehreren Jahren fremd geworden sind, ein Urteil aus dem Gedächtnis erlauben darf, so würde ich meinen, dass \textsc{Ampère}, wenn er noch lebte, von vornherein entschieden dagegen protestieren würde, wenn Sie sein Grundgesetz durch die Formel
\[
-\frac{\alpha\alpha'}{rr}ii'\sin\theta\sin\theta'\cos\varepsilon\tag*{(I)}
\]
ausdrücken. Denn das Gesetz ist etwas ganz anderes, nämlich das, was in der Formel
\[
-\frac{\alpha\alpha'}{rr}ii'\left(\tfrac12\cos\theta\cos\theta'+\sin\theta\sin\theta'\cos\varepsilon\right)\tag*{(II)}
\]
steht. Ich glaube auch nicht, dass \textsc{Ampère} mit der zusätzlichen Anmerkung zufrieden wäre, die Sie in einem späteren Brief erwähnen. Sie stellen den Unterschied dort so dar, dass \textsc{Ampères} Formel eine \textit{allgemeinere} sei, von der Form $-\dfrac{\alpha\alpha'}{rr}(F\cos\theta\cos\theta'+G\sin\theta\sin\theta'\cos\varepsilon)$. Dabei habe \textsc{Ampère} aus Versuchen $F=\tfrac12G$ abgeleitet, während Sie glauben, mit demselben Recht den Wert $F=0$ ansetzen zu können, weil seine Versuche nicht sehr genau seien. In jedem anderen als dem vorliegenden Fall würde ich zugeben, dass ein Dritter sich zu dieser Meinungsverschiedenheit zwischen Ihnen und \textsc{Ampère} etwa so äußern könnte.

\SourceLeaf{634}
Ob man dies (wie Sie) nur als eine Abwandlung des \textsc{Ampère}schen Gesetzes ansieht oder

ob es (wie \textsc{Ampère} die Sache meines Erachtens ansehen müsste) darauf hinausläuft, seine Grundformel vollständig umzustürzen und durch eine wesentlich andere zu ersetzen,

sei im Grunde doch kaum mehr als ein müßiger Streit um Worte. Wie gesagt, in jedem anderen Fall würde ich das gern einräumen, da niemand \textit{in verbis facilior} sein kann als ich. Im vorliegenden Fall ist der Unterschied aber eine Lebensfrage. Die ganze Theorie \textsc{Ampères} von der Austauschbarkeit des Magnetismus mit galvanischen Strömen hängt vollständig von der Richtigkeit der Formel II ab und geht ganz verloren, wenn an ihrer Stelle eine andere gewählt wird.

Ich kann Ihnen nicht widersprechen, wenn Sie die Versuche von \textsc{Ampère} für nicht sehr schlüssig halten, zumal ich seine klassische Abhandlung nicht zur Hand habe und mich überhaupt nicht daran erinnere, wie er seine Versuche durchgeführt hat. Doch selbst wenn er die Unvollkommenheit seiner Versuche einräumen würde, glaube ich dennoch nicht, dass er es für zulässig hielte, eine ganz andere Formel (I) zu übernehmen, an der seine ganze Theorie zerfiele, solange diese andere Formel nicht durch \textit{ganz entscheidende} Versuche bestätigt wäre. Die Bedenken, die ich Ihrem zweiten Brief zufolge selbst geäußert habe, müssen Sie missverstanden haben. Ich bin früh zu der Überzeugung gelangt und dabei geblieben, dass die oben erwähnte Austauschbarkeit \textit{notwendig} die Formel II von \textsc{Ampère} verlangt und nur andere Formeln zulässt, die für einen geschlossenen Strom mit ihr identisch werden, \textit{wenn die Wirkung in Richtung der Geraden erfolgen soll, die die beiden Stromelemente verbindet}. Lässt man diese Bedingung fallen, kann man allerdings unzählige andere Formen wählen, die aber für einen geschlossenen Strom stets dasselbe Endergebnis liefern müssen wie \textsc{Ampères} Formel. Man könnte übrigens hinzufügen, dass es für jene Zwecke stets nur um Wirkungen in messbaren Entfernungen geht. Nichts würde uns also daran hindern, anzunehmen, dass noch andere Terme zur Formel hinzukommen könnten, die nur in unmessbar kleinen Entfernungen wirksam sind (so wie zur Gravitation die molekulare Anziehung hinzutritt). Dadurch könnte die Schwierigkeit beseitigt werden, die sich aus der Abstoßung zweier aufeinanderfolgender Elemente desselben Stroms ergibt.

\SourceLeaf{635}
Um einem Missverständnis vorzubeugen, möchte ich noch darauf hinweisen, dass man die obige Formel II auch so schreiben kann
\[
-\frac{\alpha\alpha'}{rr}ii'\left(-\tfrac12\cos\theta\cos\theta'+\sin\theta\sin\theta'\cos\varepsilon\right)
\]
und dass ich nicht weiß, ob \textsc{Ampère} (dessen Abhandlung ich, wie gesagt, nicht zur Hand habe) die erste oder die zweite Schreibweise benutzt hat. Beide bedeuten dasselbe. Die erste Form schreibt man, wenn man die Winkel $\theta$, $\theta'$ an derselben (begrenzten) geraden Linie misst, diese Linie also beim zweiten Winkel in der entgegengesetzten Richtung als Schenkel wählt. Die andere Form schreibt man, wenn man eine gerade Linie unbestimmter Länge betrachtet und sie zur Messung der Winkel $\theta$, $\theta'$ beide Male in derselben Richtung heranzieht. Ebenso kann man vor die ganze Formel statt des $-$-Zeichens ein $+$-Zeichen setzen, wenn man die Anziehung anstelle der Abstoßung als positive Wirkung betrachtet.

Vielleicht gelingt es mir, mich wieder etwas mehr in diese mir inzwischen so fremd gewordenen Dinge einzuarbeiten, bis Sie mich, wie Sie mich hoffen ließen, Ende April oder Anfang Mai mit einem Besuch erfreuen. Ohne Zweifel hätte ich meine Untersuchungen längst bekannt gemacht, wenn bei ihrem Abbruch nicht gerade das gefehlt hätte, was ich als ihren eigentlichen Schlussstein ansah
\begin{center}\textit{Nil actum reputans si quid superesset agendum}\end{center}
nämlich \textit{die Ableitung} der Zusatzkräfte (die zur gegenseitigen Wirkung ruhender Teile der Elektrizität hinzukommen, wenn diese sich gegeneinander bewegen) \textit{aus der nicht augenblicklichen}, sondern sich mit der Zeit fortpflanzenden Wirkung (ähnlich wie beim Licht). Das war mir damals nicht gelungen. Soweit ich mich erinnere, ließ ich die Untersuchung aber doch nicht ganz ohne Hoffnung liegen, dass es später vielleicht gelingen könnte, wenn auch, sofern ich mich richtig erinnere, mit der persönlichen Überzeugung, dass man sich vorher eine konstruierbare Vorstellung davon machen müsse, \textit{wie} die Fortpflanzung geschieht.
\begin{center}—\quad —\quad —\quad —\quad —\quad —\quad —\quad —\quad —\quad —\quad —\quad —\quad —\quad —\quad —\quad —\end{center}
\noindent\begin{minipage}{\linewidth}
Mit herzlichen Grüßen an Ihre Geschwister und an Herrn Prof. \textsc{Möbius}

Göttingen, 19. März 1845.\hfill stets der Ihre
\begin{flushright}C. F. \textsc{Gauss}.\end{flushright}
\end{minipage}

\par\Needspace{29\baselineskip}
\SourceLeaf{636}
\begin{center}[21.]\end{center}
\noindent Lineargröße $=$ Widerstand eines gegebenen Drahts $=r$\par
\noindent Zeitgröße $=t$\par
\noindent Geschwindigkeit $=\dfrac rt$\par
\noindent Dichte $=\dfrac1{tt}=\dfrac p{r^3}$, die Dichte des Wassers beträgt etwa $\dfrac1{15\,000\,000(1'')^2}$\par
\noindent Expansibilität der Flüssigkeit $=\dfrac{rr}{t^4}=\dfrac p{rtt}$\par
\noindent Spezifische Elastizität bei einer bestimmten Temperatur $=\dfrac{rr}{tt}$\par
\noindent Beschleunigungskraft $=\dfrac r{tt}$\par
\noindent Masse $=\dfrac{r^3}{tt}=p$\par
\noindent Druck $=\dfrac{r^4}{t^4}=\dfrac{pr}{tt}$\par
\noindent Wirkung $=$ Lebend. Kraft $=$ Drehmoment $=\dfrac{r^5}{t^4}=\dfrac{prr}{tt}$\par
\noindent Wirksamkeit $=\dfrac{r^5}{t^5}=\dfrac{prr}{t^3}$\par
\noindent Erdmagnetismus $=\BareRoot{\dfrac p{rtt}}$\par
\noindent Freier Magnetismus $=$ Stärke eines ganzen Stroms $=\BareRoot{\dfrac{pr^3}{tt}}$\par
\noindent Spezifische Intensität eines galvanischen Stroms $=\BareRoot{\dfrac{pr}{tt}}$\par
\noindent Erregungskraft von Kupfer: Zink $=\BareRoot{\dfrac{pr^3}{t^4}}$\par
\noindent Leitungsvermögen eines bestimmten Metalls $=\dfrac t{rr}$\par

